\section{Task 1: Rainfall patterns}
\subsection{Across-year distribution}
The annual summaries use each retained five-minute reading as one observation.
Wet frequency is the fraction of readings greater than zero. It was lowest in
2019 (1.94\%) and highest in 2023 (3.20\%). Conditional wet-reading mean was
between 0.83 and 0.91~mm across the years. The maximum retained reading ranged
from 15.4~mm in 2019 to 37.2~mm in 2022. These annual values use all valid
records and do not adjust for changes in station coverage.
Table~\ref{tab:annual-distribution} reports the annual values, and
Figure~\ref{fig:distribution} shows the distribution and wet-reading tail.

\begin{table}[htbp]
\centering\scriptsize
\caption{Annual rainfall distribution. Mean rainfall is per retained
five-minute reading; wet mean is conditional on a positive reading.}
\label{tab:annual-distribution}
\begin{tabular}{r r r r r r}
\toprule
Year & Valid readings & Wet readings (\%) & Mean (mm) & Wet mean (mm) & Max (mm) \\
\midrule
2017 & 5,256,020 & 2.83 & 0.0235 & 0.829 & 18.0 \\
2018 & 4,594,246 & 2.88 & 0.0251 & 0.870 & 15.6 \\
2019 & 5,026,862 & 1.94 & 0.0164 & 0.847 & 15.4 \\
2020 & 6,334,473 & 3.04 & 0.0262 & 0.860 & 20.6 \\
2021 & 7,044,646 & 3.04 & 0.0277 & 0.910 & 26.4 \\
2022 & 7,133,230 & 3.09 & 0.0278 & 0.899 & 37.2 \\
2023 & 6,742,890 & 3.20 & 0.0272 & 0.849 & 17.622 \\
2024 & 6,406,260 & 3.05 & 0.0258 & 0.846 & 15.912 \\
\bottomrule
\end{tabular}
\end{table}

\begin{figure}[htbp]
\centering
\includegraphics[width=\linewidth]{\figpath/task1/rainfall_distribution.png}
\caption{Wet-reading frequency and conditional wet-reading distribution by
year. The dashed line marks the wet-reading 99th percentile estimated from
2017--2022 training-period observations; the last histogram bin contains
readings of at least 20~mm.}
\label{fig:distribution}
\end{figure}

\subsection{Within-year monthly and hourly patterns}
For eligible station-months, mean observed monthly rainfall ranged from
31.97~mm (February 2021) to 419.07~mm (October 2022). Monthly wet frequency
ranged from 0.37\% (August 2019) to 7.07\% (January 2024). The extrema do not
come from a single year, so they show variation within the observed period
rather than a stable seasonal climatology.

Conditional wet intensity also varied by month. Across the eligible
year-month summaries, it averaged 1.05~mm in March and 1.02~mm in April,
compared with 0.69~mm in January; support ranged from six to eight years per
calendar month. Individual summaries ranged from 0.58~mm (December 2019) to
1.35~mm (March 2019).

The equal station-month hourly summaries peaked at 15:00 SGT in 2018, 2019,
2021, 2022 and 2023, and at 16:00 in 2017, 2020 and 2024. The peak-hour
frequencies ranged from 3.42\% in 2019 to 7.76\% in 2023. The afternoon peak
is consistent across years, while its exact hour and magnitude vary. The
station-month extrema are listed in Table~\ref{tab:monthly-extrema};
Figures~\ref{fig:monthly} and~\ref{fig:hourly} show monthly intensity,
frequency and within-day patterns.

\begin{table}[htbp]
\centering\scriptsize
\caption{Months with the lowest and highest station-averaged wet-reading
frequency in each year. Each value averages eligible station-months with at
least 95\% monthly coverage; the support column gives the station count for
that month.}
\label{tab:monthly-extrema}
\begin{tabular}{r l r r l r}
\toprule
Year & Lowest month & Frequency (\%) & Stations & Highest month & Frequency (\%) \\
\midrule
2017 & Jun & 1.64 & 26 & Nov & 4.53 \\
2018 & Jul & 1.40 & 49 & Jan & 6.17 \\
2019 & Aug & 0.37 & 48 & Dec & 5.79 \\
2020 & Mar & 1.24 & 45 & Sep & 4.95 \\
2021 & Feb & 0.39 & 68 & Aug & 5.01 \\
2022 & Jul & 1.60 & 67 & Oct & 5.17 \\
2023 & Oct & 1.21 & 62 & Mar & 5.19 \\
2024 & Jul & 1.16 & 58 & Jan & 7.07 \\
\bottomrule
\end{tabular}
\end{table}

\begin{figure}[htbp]
\centering
\includegraphics[width=\linewidth]{\figpath/task1/monthly_frequency_intensity.png}
\caption{Monthly wet-reading frequency and conditional wet intensity by year.
Values are equal-weight averages over station-months meeting the 95\% coverage
rule; frequency uses valid five-minute readings as its denominator.}
\label{fig:monthly}
\end{figure}

\begin{figure}[htbp]
\centering
\includegraphics[width=\linewidth]{\figpath/task1/monthly_hourly_patterns.png}
\caption{Station-averaged monthly rainfall and hourly wet frequency by year.
Only station-months and station-hours meeting the coverage rules are included.
Hourly labels refer to the recorded interval-end hour in SGT.}
\label{fig:hourly}
\end{figure}

\FloatBarrier
\subsection{Network coverage and annual comparison}
Annual coverage uses each ID's full calendar year; a station-year without
readings has zero coverage. No ID met any of the 90\%, 95\% or 99\% cutoffs in
all eight years, so no fixed-ID cohort could be formed. The annual results
therefore summarize the available station network rather than a fixed
monitoring network. A screen of coordinates, station names and device IDs
found no obvious replacements or co-located IDs. The station-year coverage
pattern is shown in Figure~\ref{fig:coverage}.

\begin{figure}[!ht]
\centering
\includegraphics[width=0.80\linewidth]{\figpath/task1/coverage_heatmap.png}
\caption{Annual valid-slot coverage by station from 2017 through 2024. Blank
cells represent a station-year with no retained observations.}
\label{fig:coverage}
\end{figure}

\clearpage
\subsection{Additional descriptive views}

These figures show complete-day rainfall and coverage, hourly wet frequency,
rainfall tails, wet-spell duration and station-level summaries.
Figure~\ref{fig:complete-daily} separates the monthly rainfall level from the
completeness of the underlying observations. Monthly means on complete
station-days vary between years without a recurring month-by-month pattern,
and complete-day support also changes across the study period.

\begin{figure}[htbp]
\centering
\includegraphics[width=0.96\linewidth]{\figpath/task1/complete_daily_patterns.png}
\caption{Monthly equal-station mean rainfall on complete station-days (left)
and the fraction of expected station-days with all 288 readings (right).}
\label{fig:complete-daily}
\end{figure}

Figure~\ref{fig:hourly-supplement} shows higher wet-reading frequency in the
afternoon, generally around 15:00--17:00 SGT, than overnight. The peak hour
and its magnitude vary by year, while 2019 is lower across much of the day.

\begin{figure}[htbp]
\centering
\includegraphics[width=0.96\linewidth]{\figpath/task1/hourly_by_year.png}
\caption{Equal station-month wet-reading frequency by hour and year, using
the 95\% monthly and within-hour coverage requirements.}
\label{fig:hourly-supplement}
\end{figure}

Figure~\ref{fig:rainfall-tail} compares upper quantiles among wet readings
with the share of all valid readings above the training-period wet-reading
P99 threshold. P90 and P95 vary less than P99, while the exceedance share is
lowest in 2019, highest in 2021, and then declines through 2024.

\begin{figure}[htbp]
\centering
\includegraphics[width=0.96\linewidth]{\figpath/task1/rainfall_tail.png}
\caption{Wet-reading rainfall quantiles (P90, P95 and P99; left) and the
share of valid readings above the training-period wet-reading P99 threshold
(right).}
\label{fig:rainfall-tail}
\end{figure}

Figure~\ref{fig:wet-spells} summarizes uninterrupted sequences of positive
readings at exact five-minute intervals. The median spell is one reading in
every year, while the upper quantiles capture longer runs. Censoring at gaps
or study boundaries is more common in 2017--2018; observed spell lengths can
therefore be truncated.

\begin{figure}[htbp]
\centering
\includegraphics[width=0.96\linewidth]{\figpath/task1/wet_spell_patterns.png}
\caption{Consecutive wet-reading spells at individual stations. Missing
intervals censor spells; the figure is supplementary and does not imply
continuous rainfall between observations.}
\label{fig:wet-spells}
\end{figure}

Figure~\ref{fig:station-spatial} maps station-year summaries with at least
90\% annual coverage and consistent coordinates through 2023. Point colour
shows observed annual rainfall. Totals vary across nearby locations, but
these station observations do not estimate rainfall between stations or an
area-weighted spatial pattern.

\begin{figure}[htbp]
\centering
\includegraphics[width=0.82\linewidth]{\figpath/task1/station_spatial_patterns.png}
\caption{Observed annual rainfall at station-years with at least 90\%
coverage and consistent coordinates through 2023. Each point is a station-year;
the plot is not an area-weighted map of Singapore rainfall.}
\label{fig:station-spatial}
\end{figure}
\clearpage

